\chapter{Forking lemmas}

The forking lemma (Pointcheval--Stern, Bellare--Neven) is the central rewinding
argument behind the security of Fiat--Shamir signatures in the random-oracle
model. Running an adversary twice on the same commitment but with fresh oracle
responses yields two accepting transcripts, from which special-soundness
extraction recovers a witness. This chapter records the single-query bound and
its Bellare--Neven generalisation to $q$ oracle queries.

\section{The forkable adversary}

\begin{definition}[Forkable adversary]
  \label{def:forkable-adversary}
  \lean{CatCrypt.Crypto.ForkingLemma.ForkableAdversary}
  \leanok
  A \emph{forkable adversary} for the single-query random-oracle model is a
  two-phase machine over input $X$, commitment $C$, response space $R$, and
  output $Y$: a probabilistic commit phase
  $\mathsf{coins} : X \to \mathsf{SPComp}\,C$ (everything computed before the
  oracle query) and a deterministic response phase
  $\mathsf{respond} : X \to C \to R \to Y$ (the computation after the oracle
  answers). Sharing the commit phase across two runs is what enables forking.
\end{definition}

\begin{definition}[Acceptance probability]
  \label{def:fork-accept-prob}
  \uses{def:forkable-adversary}
  \lean{CatCrypt.Crypto.ForkingLemma.acceptProb}
  \leanok
  The \emph{acceptance probability} of $A$ on input $x$ under predicate
  $\mathsf{accept} : Y \to \mathsf{Bool}$ is
  \[
    \mathrm{acc}(A,x) \;=\;
    \Pr\big[\, c \leftarrow \mathsf{coins}\,x;\;
      e \leftarrow_{\$} R \;:\;
      \mathsf{accept}(\mathsf{respond}\,x\,c\,e) \,\big].
  \]
\end{definition}

\begin{definition}[Forking algorithm]
  \label{def:fork-alg}
  \uses{def:forkable-adversary}
  \lean{CatCrypt.Crypto.ForkingLemma.forkAlg}
  \leanok
  The \emph{forking algorithm} runs $A$ once to obtain a commitment $c$, then
  samples two independent oracle responses $e_1, e_2 \leftarrow_{\$} R$ and
  replays the response phase on each, returning
  $(c, e_1, y_1, e_2, y_2)$ with $y_i = \mathsf{respond}\,x\,c\,e_i$.
\end{definition}

\begin{definition}[Fork success probability]
  \label{def:fork-succ-prob}
  \uses{def:fork-alg}
  \lean{CatCrypt.Crypto.ForkingLemma.forkSuccProb}
  \leanok
  The \emph{fork success probability} is the probability that both replayed runs
  accept on distinct oracle responses:
  \[
    \mathrm{frk}(A,x) \;=\;
    \Pr\big[\,\mathsf{accept}\,y_1 \wedge \mathsf{accept}\,y_2 \wedge e_1 \neq e_2\,\big],
  \]
  taken over the forking algorithm.
\end{definition}

\section{Jensen's inequality for sub-PMFs}

\begin{lemma}[AM--GM for extended non-negative reals]
  \label{lem:ennreal-amgm}
  \lean{CatCrypt.Crypto.ForkingLemma.ennreal_two_mul_le_add_sq}
  \leanok
  For all $a, b \in \overline{\mathbb{R}}_{\geq 0}$,
  \[
    2ab \;\leq\; a^2 + b^2 .
  \]
\end{lemma}

\begin{lemma}[Jensen's inequality under sub-PMF weights]
  \label{lem:jensen-sub-pmf}
  \uses{lem:ennreal-amgm}
  \lean{CatCrypt.Crypto.ForkingLemma.jensen_sub_pmf_sq}
  \leanok
  Let $w : \iota \to \overline{\mathbb{R}}_{\geq 0}$ with
  $\sum_i w_i \leq 1$ and $f : \iota \to \overline{\mathbb{R}}_{\geq 0}$ with
  $f_i \leq 1$ for all $i$. Then squaring is convex against the weights $w$:
  \[
    \Big(\sum_i w_i\, f_i\Big)^2 \;\leq\; \sum_i w_i\, f_i^{\,2} .
  \]
\end{lemma}

\section{The conditional fork, per commitment}

\begin{definition}[Conditional acceptance probability]
  \label{def:cond-acc-prob}
  \uses{def:forkable-adversary}
  \lean{CatCrypt.Crypto.ForkingLemma.condAccProb}
  \leanok
  Fixing a commitment $c$, the \emph{conditional acceptance probability}
  $\mathrm{acc}_c$ is the probability that a single freshly sampled oracle
  response $e \leftarrow_{\$} R$ makes $\mathsf{respond}\,x\,c\,e$ accept.
\end{definition}

\begin{definition}[Conditional fork probability]
  \label{def:cond-fork-prob}
  \uses{def:forkable-adversary}
  \lean{CatCrypt.Crypto.ForkingLemma.condForkProb}
  \leanok
  Fixing a commitment $c$, the \emph{conditional fork probability}
  $\mathrm{frk}_c$ is the probability that two independent responses
  $e_1, e_2 \leftarrow_{\$} R$ both accept and are distinct.
\end{definition}

\begin{lemma}[Exact conditional fork count]
  \label{lem:cond-fork-prob-eq}
  \uses{def:cond-acc-prob, def:cond-fork-prob}
  \lean{CatCrypt.Crypto.ForkingLemma.condForkProb_eq}
  \leanok
  For each fixed commitment $c$, exact counting over $R \times R$ gives
  \[
    \mathrm{frk}_c \;=\; \mathrm{acc}_c^{\,2} \;-\; \frac{\mathrm{acc}_c}{\lvert R\rvert},
  \]
  the off-diagonal accepting pairs divided by $\lvert R\rvert^2$.
\end{lemma}

\begin{lemma}[Decomposition over commitments]
  \label{lem:fork-decompose}
  \uses{def:fork-accept-prob, def:fork-succ-prob, def:cond-acc-prob, def:cond-fork-prob}
  \lean{CatCrypt.Crypto.ForkingLemma.acceptProb_decompose, CatCrypt.Crypto.ForkingLemma.forkSuccProb_decompose}
  \leanok
  Both $\mathrm{acc}(A,x)$ and $\mathrm{frk}(A,x)$ decompose as weighted sums over
  the commitment PMF: $\mathrm{acc} = \sum_c w_c\,\mathrm{acc}_c$ and
  $\mathrm{frk} = \sum_c w_c\,\mathrm{frk}_c$, where $w_c$ is the probability the
  commit phase outputs $c$.
\end{lemma}

\section{The single-query forking lemma}

\begin{theorem}[Forking lemma, $q = 1$]
  \label{thm:fork-success-ge}
  \uses{lem:jensen-sub-pmf, lem:cond-fork-prob-eq, lem:fork-decompose}
  \lean{CatCrypt.Crypto.ForkingLemma.fork_success_ge}
  \leanok
  For every forkable adversary $A$, input $x$, and acceptance predicate,
  \[
    \mathrm{frk}(A,x) \;\geq\;
    \mathrm{acc}(A,x)^2 \;-\; \frac{\mathrm{acc}(A,x)}{\lvert R\rvert}.
  \]
  The bound follows by decomposing over commitments, applying the exact
  conditional identity per commitment, and using Jensen's inequality to move the
  square outside the commitment average.
\end{theorem}

\section{The Bellare--Neven multi-query lemma}

\begin{definition}[$q$-query forkable adversary]
  \label{def:qforkable-adversary}
  \lean{CatCrypt.Crypto.GeneralForkingLemma.QForkableAdversary}
  \leanok
  A \emph{$q$-query forkable adversary} has a probabilistic commit phase
  $\mathsf{coins} : X \to \mathsf{SPComp}\,C$ and a deterministic response phase
  $\mathsf{respond} : X \to C \to (\mathrm{Fin}\,q \to R) \to \mathrm{Fin}\,q \times Y$
  that, given all $q$ oracle responses, outputs an index identifying the relevant
  query together with its result.
\end{definition}

\begin{definition}[$q$-query acceptance probability]
  \label{def:qaccept-prob}
  \uses{def:qforkable-adversary}
  \lean{CatCrypt.Crypto.GeneralForkingLemma.qAcceptProb}
  \leanok
  $\mathrm{acc}(A,x)$ for a $q$-query adversary is the probability that, over a
  commitment and a uniform response vector $\mathrm{Fin}\,q \to R$, the predicate
  accepts the produced $(\text{index}, \text{output})$ pair.
\end{definition}

\begin{definition}[Guess reduction to a single query]
  \label{def:guess-reduction}
  \uses{def:qforkable-adversary, def:forkable-adversary}
  \lean{CatCrypt.Crypto.GeneralForkingLemma.guessReduction}
  \leanok
  For a guess $I : \mathrm{Fin}\,q$, the \emph{guess reduction} builds a
  single-query forkable adversary: its commit phase absorbs the commitment and a
  full uniform response vector, and its response phase overwrites position $I$
  with the single oracle challenge via $\mathsf{Function.update}$ before running
  $\mathsf{respond}$.
\end{definition}

\begin{definition}[Acceptance at a guess]
  \label{def:accept-at-guess}
  \lean{CatCrypt.Crypto.GeneralForkingLemma.acceptAtGuess}
  \leanok
  The predicate $\mathsf{acceptAtGuess}(\mathsf{accept}, I)$ accepts a pair
  $(j, y)$ exactly when $j = I$ and $\mathsf{accept}(j,y)$ holds; it restricts
  attention to transcripts whose fork index matches the guess.
\end{definition}

\begin{lemma}[Guesses partition acceptance]
  \label{lem:sum-accept-at-guess}
  \uses{def:accept-at-guess, def:guess-reduction, def:qaccept-prob}
  \lean{CatCrypt.Crypto.GeneralForkingLemma.sum_acceptAtGuess_eq}
  \leanok
  Summing the reduced acceptance probabilities over all guesses recovers the
  $q$-query acceptance probability:
  \[
    \sum_{I : \mathrm{Fin}\,q}
      \mathrm{acc}\big(\mathsf{guessReduction}(A,I),\, x,\, \mathsf{acceptAtGuess}(\mathsf{accept},I)\big)
    \;=\; \mathrm{acc}(A,x).
  \]
\end{lemma}

\begin{theorem}[Bellare--Neven aggregation]
  \label{thm:bellare-neven-agg}
  \uses{lem:ennreal-amgm}
  \lean{CatCrypt.Crypto.GeneralForkingLemma.bellare_neven_aggregation}
  \leanok
  Let $q > 0$ and $a, f : \mathrm{Fin}\,q \to \overline{\mathbb{R}}_{\geq 0}$ with
  $f_j \geq a_j^2 - a_j / N$ for every $j$. Then, by Cauchy--Schwarz on the
  per-index bounds,
  \[
    \sum_{j} f_j \;\geq\;
    \frac{\big(\sum_j a_j\big)^2}{q} \;-\; \frac{\sum_j a_j}{N}.
  \]
\end{theorem}

\begin{definition}[Indexed $q$-fork success probability]
  \label{def:qfork-succ-prob}
  \uses{def:guess-reduction, def:accept-at-guess, def:fork-succ-prob}
  \lean{CatCrypt.Crypto.GeneralForkingLemma.qForkSuccProb}
  \leanok
  The \emph{indexed $q$-fork success probability} sums the single-query fork
  success of every guess reduction, each measured against its guess-restricted
  predicate:
  \[
    \mathrm{qfrk}(A,x) \;=\;
    \sum_{I : \mathrm{Fin}\,q}
      \mathrm{frk}\big(\mathsf{guessReduction}(A,I),\, x,\, \mathsf{acceptAtGuess}(\mathsf{accept},I)\big).
  \]
\end{definition}

\begin{theorem}[General forking lemma, Bellare--Neven 2006]
  \label{thm:qfork-success-ge}
  \uses{thm:fork-success-ge, thm:bellare-neven-agg, lem:sum-accept-at-guess, def:qfork-succ-prob, def:qaccept-prob}
  \lean{CatCrypt.Crypto.GeneralForkingLemma.qFork_success_ge}
  \leanok
  For every $q$-query forkable adversary $A$ with $q > 0$,
  \[
    \mathrm{qfrk}(A,x) \;\geq\;
    \frac{\mathrm{acc}(A,x)^2}{q} \;-\; \frac{\mathrm{acc}(A,x)}{\lvert R\rvert}.
  \]
  Combining the single-query bound applied to each guess reduction, the guess
  partition of acceptance, and the Cauchy--Schwarz aggregation yields the
  Bellare--Neven bound.
\end{theorem}
